\documentclass[12pt]{article}

\usepackage{amsmath}
\usepackage{amssymb}
\usepackage[paper=a4paper]{geometry}
% \usepackage{caption}

\author{Zac Hart}
\title{\bf Extension 1 - Tutorial 1}

\begin{document}

\maketitle

\section*{Answers}

\begin{enumerate}
    % 1.
    \item {\bf Prove by mathematical induction that
    \begin{equation*}
        \mathbf{\frac{1}{1 \times 3} + \frac{1}{3 \times 5} + \frac{1}{5 \times 7} + \dots + \frac{1}{(2n - 1)(2n + 1)} = \frac{n}{2n + 1}}
    \end{equation*}
    for all positive integer values of $\mathbf{n}$.} \\

    In the case where $n = 1$,
    \begin{flalign*}
        && \text{RHS} &= \frac{1}{2(1)+1} = \frac{1}{3}, & \\
        \text{and} && \text{LHS} &= \frac{1}{1 \times 3} = \frac{1}{3} = \text{RHS}.
    \end{flalign*}
    The result is thus true for the initial case.

    Now assume that the result is true for the positive integer $n = k$, that is,
    \begin{equation}
        \label{eq:1}
        \frac{1}{1 \times 3} + \frac{1}{3 \times 5} + \frac{1}{5 \times 7} + \cdots + \frac{1}{(2k - 1)(2k + 1)} = \frac{k}{2k +1}.
    \end{equation}

    We want to show that the result is true for $n = k + 1$. We have
    \begin{gather*}
        \text{RHS} = \frac{k + 1}{2k + 3}, \\
        \text{LHS} = \frac{1}{1 \times 3} + \frac{1}{3 \times 5} + \cdots + \frac{1}{(2k - 1)(2k + 1)} + \frac{1}{(2k + 1)(2k + 3)};
    \end{gather*}
    by substituting (\ref{eq:1}) into the LHS, we get
    \begin{align*}
        \text{LHS} &= \frac{k}{2k + 1} + \frac{1}{(2k + 1)(2k + 3)} \\
        &= \frac{k(2k + 3) + 1}{(2k + 1)(2k + 3)} \\
        &= \frac{2k^2 + 3k + 1}{(2k + 1)(2k + 3)} \\
        &= \frac{(2k + 1)(k + 1)}{(2k + 1)(2k + 3)} \\
        &= \frac{(k + 1)}{(2k + 3)} \\
        &= \text{RHS}.
    \end{align*}

    The result is thus true for $n = k + 1$. But it is true for $n = 1$, hence it is also true for $n = 2, 3, \dots$, extending to all positive integer values of $n$.

    % 2.
    \item {\bf Prove by mathematical induction that
    \begin{equation*}
        \mathbf{1^3 + 2^3 + 3^3 + \dots + n^3 = \frac{1}{4}n^2(n + 1)^2}
    \end{equation*}
    for all positive integer values of $\mathbf{n}$.} \\

    In the case where $n = 1$,
    \begin{flalign*}
        && \text{RHS} &= \frac{1}{4}(1)^2(1 + 1)^2 & \\
        && &= \frac{4}{4} = 1, \\
        \text{and} && \text{LHS} &= 1^3 = 1 \\
        && &= \text{RHS}.
    \end{flalign*}
    The result is thus true for the initial case.

    Assume that the result is true for the positive integer $n = k$, that is,
    \begin{equation}
        \label{eq:2}
        1^3 + 2^3 + \dots + k^3 = \frac{1}{4}k^2(k + 1)^2.
    \end{equation}

    When $n = k + 1$,
    \begin{align*}
        \text{RHS} &= \frac{1}{4}(k + 1)^2(k + 2)^2, \\
        \text{LHS} &= 1^3 + 2^3 + \dots + k^3 + (k + 1)^3;
    \end{align*}
    by substituting (\ref{eq:2}) into the LHS, we get
    \begin{align*}
        \text{LHS} &= \frac{1}{4}k^2(k + 1)^2 + (k + 1)^3 \\
        &= \frac{1}{4}(k + 1)^2(k^2 + 4k + 4) \\
        &= \frac{1}{4}(k + 1)^2(k + 2)^2 \\
        &= \text{RHS}.
    \end{align*}

    The result is thus true for $n = k + 1$. But it is true for $n = 1$, hence it is also true for all positive integer values of $n$.

    % 3.
    \item {\bf Prove by mathematical induction that
    \begin{equation*}
        \mathbf{\sum_{r = 1}^n \frac{r}{r^4 + r^2 + 1} = \frac{n(n + 1)}{2(n^2 + n + 1)}}
    \end{equation*}
    for all positive integer values of $\mathbf{n}$.} \\

    In the case where $n = 1$,
    \begin{flalign*}
        && \text{RHS} &= \frac{1(1 + 1)}{2(1^2 + 1 + 1)} & \\
        && &= \frac{2}{2(3)} = \frac{1}{3}, \\
        \text{and} && \text{LHS} &= \frac{1}{1^4 + 1^2 + 1} \\
        && &= \frac{1}{3} \\
        && &= \text{RHS}.
    \end{flalign*}
    The result is thus true for the initial case.

    Assume that the result is true for the positive integer $n = k$, that is,
    \begin{equation}
        \label{eq:3}
        \sum_{r = 1}^k \frac{r}{r^4 + r^2 + 1} = \frac{k(k + 1)}{2(k^2 + k + 1)}.
    \end{equation}

    When $n = k + 1$,
    \begin{flalign*}
        && \text{RHS} &= \frac{(k + 1)((k + 1) + 1)}{2((k + 1)^2 + (k + 1) + 1)} & \\
        && &= \frac{(k + 1)(k + 2)}{2(k^2 + 3k + 3)}, \\
        \text{and} && \text{LHS} &= \sum_{r = 1}^{k + 1} \frac{r}{r^4 + r^2 + 1} \\
        && &= \sum_{r = 1}^k \frac{r}{r^4 + r^2 + 1} + \frac{k + 1}{(k + 1)^4 + (k + 1)^2 + 1} \\
        && &= \frac{k(k + 1)}{2(k^2 + k + 1)} + \frac{k + 1}{(k + 1)^4 + (k + 1)^2 + 1}. & \text{[from (\ref{eq:3})]}
    \end{flalign*}

    In order to simplify the $LHS$, we will first factorise $(k + 1)^4 + (k + 1)^2 + 1$.

    Since $n = k + 1$, we have
    \begin{align*}
        (k + 1)^4 + (k + 1)^2 + 1 &= n^4 + n^2 + 1.
        \intertext{Completing the square,}
        (k + 1)^4 + (k + 1)^2 + 1 &= (n^4 + 2n^2 + 1) - n^2 \\
        &= (n^2 + 1)^2 - n^2 \\
        &= (n^2 + n + 1)(n^2 - n + 1),
        \intertext{and since $n = k + 1$,}
        (k + 1)^4 + (k + 1)^2 + 1 &= ((k + 1)^2 + (k + 1) + 1)((k + 1)^2 - (k + 1) + 1) \\
        &= (k^2 + 3k + 3)(k^2 + k + 1).
    \end{align*}
    The $LHS$ is thus
    \begin{align*}
        \text{LHS} &= \frac{k(k + 1)}{2(k^2 + k + 1)} + \frac{k + 1}{(k^2 + 3k + 3)(k^2 + k + 1)} \\
        &= \frac{k(k + 1)(k^2 + 3k + 3) + 2(k + 1)}{2(k^2 + 3k + 3)(k^2 + k + 1)} \\
        &= \frac{(k + 1)(k(k^2 + 3k + 3) + 2)}{2(k^2 + 3k + 3)(k^2 + k + 1)} \\
        &= \frac{(k + 1)(k^3 + 3k^2 + 3k + 2)}{2(k^2 + 3k + 3)(k^2 + k + 1)}.
    \end{align*}

    Now we will factorise the polynomial $f(k) = k^3 + 3k^2 + 3k + 2$. Notice that
    \begin{align*}
        f(-2) &= (-2)^3 + 3(-2)^2 + 3(-2) + 2 \\
        &= -8 + 12 - 6 + 2 \\
        &= 0.
    \end{align*}
    Hence, by the remainder theorem, $(k + 2)$ is a factor of $f(k)$. Dividing $f(k)$ by $(k + 2)$, we obtain the other factor, which is $k^2 + k + 1$. Therefore,
    \begin{equation*}
        k^3 + 3k^2 + 3k + 2 = (k + 2)(k^2 + k + 1).
    \end{equation*}

    Substituting this into the $\text{LHS}$, we get
    \begin{align*}
        \text{LHS} &= \frac{(k + 1)(k + 2)(k^2 + k + 1)}{2(k^2 + 3k + 3)(k^2 + k + 1)} \\
        &= \frac{(k + 1)(k + 2)}{2(k^2 + 3k + 3)} \\
        &= \text{RHS}.
    \end{align*}

    The result is thus true for $n = k + 1$. But it is true for $n = 1$, hence it is also true for all positive integer values of $n$.

    % 4.
    \item {\bf Let $\mathbf{x^{(r)} = x(x - 1)(x - 2)(x - 3)\cdots(x - r + 1)}$, where $\mathbf{x}$ and $\mathbf{r}$ are both integers. Prove by mathematical induction that
    \begin{equation*}
        \mathbf{\sum_{x = 1}^n x^{(r)} = \frac{(n + 1)^{(r + 1)}}{r + 1},}
    \end{equation*}
    where $\mathbf{n}$ is a positive integer.} \\

    It is important to note the nature of the function $x^{(r)}$.

    Since the ordered terms in the product are decreasing from left to right, the final term must be less than the first. It is therefore necessary to restrict $r$ to the positive integers in order to satisfy this requirement.

    Additionally, the product is zero when $r \ge x$. Consider, for example, when $x = 2$ and $r = 4$. The value of $x^{(n)}$ is
    \begin{align*}
        2^{(4)} &= 2(2 - 1)(2 - 2)(2 - 3) \\
        &= 2(1)(0)(-1) \\
        &= 0.
    \end{align*}
    From this example, it is observed that there exists a term $(x - x) = 0$ in the product. This only occurs when $r \ge x$.

    The summation $\sum_{x = 1}^n x^{(r)}$ will therefore contain a number of zero terms -- except when $r = 1$, in which case $x^{(r)}$ is effectively equivalent to the factorial function.

    When $r > 1$, the greatest value of $x$ such that $x^{(r)}$ equals zero occurs when the last term in the product equals zero, that is, $(x - r) = 0$ or $x = r$. We can therefore assume that every previous term also equals zero, so that
    \begin{equation*}
        \sum_{x = 1}^n x^{(r)} = \sum_{x = r + 1}^n x^{(r)}.
    \end{equation*}

    We will now move only to the proof by induction.

    In the case where $n = 1$,
    \begin{align*}
        \text{RHS} &= \frac{(1 + 1)^{(r + 1)}}{r + 1} \\
        &= \frac{2^{(r + 1)}}{r + 1} \\
        &= \begin{cases}
            1 & \text{if } r = 1 \\
            0 & \text{if } r >1
        \end{cases},
        \intertext{and}
        \text{LHS} &= 1^{(r)} \\
        &= \begin{cases}
            1 & \text{if } r = 1 \\
            0 & \text{if } r > 1
        \end{cases} \\
        &= \text{RHS}.
    \end{align*}
    The result is thus true for the initial case.

    Assume that the result is true for the positive integer $n = k$, that is,
    \begin{equation}
        \label{eq:4}
        \sum_{x = 1}^k x^{(r)} = \sum_{x = r + 1}^k x^{(r)} = \frac{(k + 1)^{(r + 1)}}{r + 1}.
    \end{equation}

    When $n = k + 1$,
    \begin{flalign*}
        %
    \end{flalign*}

    The result is thus true for $n = k + 1$. But it is true for $n = 1$, hence it is also true for all positive integer values of $n$.

    [NOTE TO MR. HORSLEY]: I have left this question unfinished, however I will send you a new copy of this document upon completing it.

    % 5.
    \item {\bf Prove by mathematical induction that
    \begin{equation*}
        \mathbf{\sum_{k = 1}^n k(3k - 1) = n^2(n + 1)}
    \end{equation*}
    for all positive integer values of $\mathbf{n}$.} \\

    In the case where $n = 1$,
    \begin{flalign*}
        && \text{RHS} &= 1^2(1 + 1) = 2, & \\
        \text{and} && \text{LHS} &= 1(3(1) - 1) = 2 = \text{RHS}.
    \end{flalign*}
    The result is thus true for the initial case.

    Assume that the result is true for the positive integer $n = k$, that is,
    \begin{equation}
        \label{eq:5}
        \sum_{r = 1}^k r(3r - 1) = k^2(k + 1).
    \end{equation}

    When $n = k + 1$,
    \begin{flalign*}
        && \text{RHS} &= \sum_{r = 1}^k r(3r - 1) + (k + 1)(3(k + 1) - 1) \\
        && &= \sum_{r = 1}^k r(3r - 1) + (k + 1)(3k + 2) \\
        && &= k^2(k + 1) + (k + 1)(3k + 2) & \text{[from (\ref{eq:5})]} \\
        && &= (k + 1)(k^2 + 3k + 2) \\
        && &= (k + 1)^2(k + 2), \\
        \text{and} && \text{LHS} &= (k + 1)^2((k + 1) + 1) \\
        && &= (k + 1)^2(k + 2) \\
        && &= \text{RHS}.
    \end{flalign*}

    The result is thus true for $n = k + 1$. But it is true for $n = 1$, hence it is also true for all positive integer values of $n$.

    % 6.
    \item {\bf Prove by mathematical induction that $\mathbf{5^n + 2 \times 11^n}$ is divisible by 3 for all values of $\mathbf{n \ge 0}$.} \\

    In the case where $n = 0$,
    \begin{equation*}
        5^0 + 2 \times 11^0 = 1 + 2 \times 1 = 3,
    \end{equation*}
    which is a multiple of 3. The result is thus true for the initial case.

    Assume that the result is true for the positive integer $n = k$, that is,
    \begin{equation}
        \label{eq:6}
        5^k + 2 \times 11^k = 3M,
    \end{equation}
    where $M$ is an integer.

    When $n = k + 1$,
    \begin{align*}
        5^{k + 1} + 2 \times 11^{k + 1} &= 5 \times 5^k + 22 \times 11^k \\
        &= 5(5^k + 211^k) + 12 \times 11^k \\
        &= 5(3M) + 12 \times 11^k & \text{[from (\ref{eq:6})]} \\
        &= 15M + 12 \times 11^k \\
        &= 3(5M + 4 \times 11^k).
    \end{align*}
    Since $(5M + 4 \times 11^k)$ is an integer, the result is true for $n = k + 1$. But it is true for $n = 0$, hence it is also true for all $n \ge 0$.

    % 7.
    \item {\bf Consider the expression
    \begin{equation*}
        \mathbf{A^{n + 1} - A^nB + B^{n + 1} - B^nA,}
    \end{equation*}
    where $\mathbf{A}$ and $\mathbf{B}$ are positive real numbers and $\mathbf{n}$ is a positive integer.}

    \begin{enumerate}
        % 7.a.
        \item {\bf Show that $\mathbf{A^{n + 1} + B^{n + 1} \ge A^nB + B^nA}$.} \\

        We factorise the given expression to get
        \begin{align*}
            A^{n + 1} - A^nB + B^{n + 1} - B^nA &= A^n(A - B) + B^n(B - A) \\
            &= (A^n - B^n)(A - B) \\
            &= (A - B)^2(A^{n - 1} + A^{n - 2}B + \cdots + AB^{n - 2} + B^{n - 1}).
        \end{align*}
        Since both of the factors are always zero or positive,
        \begin{flalign*}
            && A^{n + 1} - A^nB + B^{n + 1} - B^nA &\ge 0, & \\
            \text{or} && A^{n + 1} + B^{n + 1} &\ge A^nB + B^nA.
        \end{flalign*}

        % 7.b.
        \item {\bf Using the result from part (a), prove by mathematical induction that
        \begin{equation*}
            \mathbf{\left(\frac{A + B}{2}\right)^n \le \frac{A^n + B^n}{2}}
        \end{equation*}
        for all positive integer values of $\mathbf{n}$.} \\

        In the case where $n = 1$,
        \begin{flalign*}
            && \text{RHS} &= \left(\frac{A + B}{2}\right)^1 = \frac{A + B}{2}, & \\
            \text{and} && \text{LHS} &= \frac{A^1 + B^1}{2} = \frac{A + B}{2} = \text{RHS}.
        \end{flalign*}
        The result is thus true for the initial case.

        Assume that the result is true for the positive integer $n = k$, that is,
        \begin{equation}
            \label{eq:7}
            \left(\frac{A + B}{2}\right)^k \le \frac{A^k + B^k}{2}.
        \end{equation}

        When $n = k + 1$,
        \begin{gather*}
        \begin{flalign*}
            && \text{RHS} &= \frac{A^{k + 1} + B^{k + 1}}{2}, & \\
            \text{and} && \text{LHS} &= \left(\frac{A + B}{2}\right)^{k + 1}.
        \end{flalign*}
        \\
        \begin{flalign*}
            \text{Hence} && \text{LHS} - \text{RHS} &= \left(\frac{A + B}{2}\right)^{k + 1} - \frac{A^{k + 1} + B^{k + 1}}{2} \\
            && &\le \left(\frac{A + B}{2}\right)\left(\frac{A^k + B^k}{2}\right) - \frac{A^{k + 1} + B^{k + 1}}{2} & \text{[from (\ref{eq:7})]} \\
            && &= \frac{A^{k + 1} + B^{k + 1} + A^kB + B^kA}{4} - \frac{A^{k + 1} + B^{k + 1}}{2} \\
            && &= \frac{A^kB + B^kA - A^{k + 1} - B^{k + 1}}{4}.
        \end{flalign*}
        \end{gather*}

        But it was previously proven that $A^{n + 1} + B^{n + 1} \ge A^nB + B^nA$, so
        \begin{flalign*}
            && \text{LHS} - \text{RHS} &= \frac{(A^kB + B^kA) - (A^{k + 1} - B^{k + 1})}{4} & \\
            && &\le \frac{(A^{k + 1} - B^{k + 1}) - (A^{k + 1} - B^{k + 1})}{4} \\
            && &= 0, \\
            \text{hence} && \text{LHS} &\le \text{RHS}.
        \end{flalign*}

        The result is thus true for $n = k + 1$. But it is true for $n = 1$, hence it is also true for all positive integer values of $n$.
    \end{enumerate}

    % 8.
    \item {\bf Prove by mathematical induction that $\mathbf{15^n + 16^{2n + 1}}$ is divisible by $\mathbf{7}$ for all $\mathbf{n \ge 0}$.} \\

    In the case where $n = 0$,
    \begin{equation*}
        15^0 + 16^{2(0) + 1} = 1 + 16^1 = 17,
    \end{equation*}
    which is a prime number. This is a counterexample of the result, and thus disproves it.

    % 9.
    \item {\bf Prove by mathematical induction that
    \begin{equation*}
        \mathbf{(1 \times 2^1) + (2 \times 2^2) + (3 \times 2^3) + \dots + (n \times 2^n) = 2 + (n - 1)2^{n + 1}}
    \end{equation*}
    for all positive integer values of $\mathbf{n}$.} \\

    In the case where $n = 1$,
    \begin{flalign*}
        && \text{RHS} &= 2 + (1 - 1)2^{1 + 1} = 2, & \\
        \text{and} && \text{LHS} &= 1 \times 2^1 = 2 = \text{RHS}.
    \end{flalign*}
    The result is thus true for the initial case.

    Assume that the result is true for the positive integer $n = k$, that is,
    \begin{equation}
        \label{eq:8}
        (1 \times 2^1) + (2 \times 2^2) + (3 \times 2^3) + \dots + (k \times 2^k) = 2 + (k - 1)2^{k + 1}.
    \end{equation}

    When $n = k + 1$,
    \begin{gather*}
    \begin{align*}
        \text{RHS} &= 2 + ((k + 1) - 1)2^{(k + 1) + 1} & \\
        &= 2 + k \times 2^{k + 2},
    \end{align*}
    \\
    \begin{flalign*}
        \text{and} && \text{LHS} &= (1 \times 2^1) + (2 \times 2^2) + \dots + (k \times 2^k) + ((k + 1) \times 2^{k + 1}) \\
        && &= (2 + (k - 1)2^{k + 1}) + (k + 1)2^{k + 1} & \text{[from (\ref{eq:8})]} \\
        && &= 2 + 2k \times 2^{k + 1} \\
        && &= 2 + k \times 2^{k + 2} \\
        && &= \text{RHS}.
    \end{flalign*}
    \end{gather*}

    The result is thus true for $n = k + 1$. But it is true for $n = 1$, hence it is also true for all positive integer values of $n$.
\end{enumerate}

\end{document}
